\documentclass[preview]{standalone}

\usepackage{amsmath}
\usepackage{amssymb}
\usepackage{stellar}
\usepackage{definitions}
\usepackage{bettelini}

\begin{document}

\id{linearalgebra-elementary-matrices}
\genpage

\section{Elementary matrices}

\includesnpt{linearalgebra-elementary-row-operations}

\begin{snippetdefinition}{elementary-lambda-matrix-definition}{Elementary matrix}
    Let \(\mathbb{F}\) be a \field, let \(n\in\naturalnumbers\), let
    \(1\leq i,j\leq n\) with \(i\neq j\), and let
    \(\lambda\in\mathbb{F}\). The \emph{elementary matrix}
    \(E_{ij}(\lambda)\in\matrices_n(\mathbb{F})\) is the \matrix that
    coincides with \(I_n\) except possibly in the \((i,j)\)-entry, where
    its value is \(\lambda\). Equivalently,
    \[
        \left(E_{ij}(\lambda)\right)_{ab}
        =
        \kronecker_{ab}
        +
        \lambda\kronecker_{ia}\kronecker_{bj}.
    \]
\end{snippetdefinition}

\begin{snippetexample}{elementary-matrix-3-by-3-example}{An elementary matrix}
    Let \(\mathbb{F}\) be a \field and let \(\lambda\in\mathbb{F}\). In
    \(\matrices_{3\times3}(\mathbb{F})\),
    \[
        E_{23}(\lambda)
        =
        \begin{pmatrix}
            1 & 0 & 0\\
            0 & 1 & \lambda\\
            0 & 0 & 1
        \end{pmatrix}.
    \]
\end{snippetexample}

\begin{snippetproposition}{elementary-matrix-row-operation}{Elementary matrices and row addition}
    Let \(\mathbb{F}\) be a \field, let
    \(A\in\matrices_{n\times m}(\mathbb{F})\), let
    \(b\in\mathbb{F}^n\), let \(i\neq j\), and let
    \(\lambda\in\mathbb{F}\). The row operation
    \[
        R_i\leftarrow R_i-\lambda R_j
    \]
    on the linear system \(Ax=b\) is obtained by left multiplying both sides
    by \(E_{ij}(-\lambda)\):
    \[
        E_{ij}(-\lambda)Ax=E_{ij}(-\lambda)b.
    \]
\end{snippetproposition}

\begin{snippetproof}{elementary-matrix-row-operation-proof}{elementary-matrix-row-operation}{Elementary matrices and row addition}
    Let \(R_k(A)\) denote the \(k\)-th row of \(A\). In the product
    \(E_{ij}(-\lambda)A\), every row except the \(i\)-th row is unchanged,
    while the \(i\)-th row is
    \[
        R_i(A)-\lambda R_j(A).
    \]
    Indeed, the only off-diagonal nonzero entry of \(E_{ij}(-\lambda)\) is
    \(-\lambda\) in position \((i,j)\). Therefore left multiplication by
    \(E_{ij}(-\lambda)\) performs exactly
    \(R_i\leftarrow R_i-\lambda R_j\) on \(A\). The same computation applied
    to the column vector \(b\) replaces \(b_i\) by \(b_i-\lambda b_j\) and
    leaves every other component unchanged. Hence the transformed system is
    \(E_{ij}(-\lambda)Ax=E_{ij}(-\lambda)b\).
\end{snippetproof}

\begin{snippetnote}{row-operations-as-left-multiplication-note}{Row operations as left multiplication}
    Row operations on a linear system correspond to left multiplication of
    both the coefficient \matrix and the constant vector by suitable
    elementary \matrix[matrices].
\end{snippetnote}

\begin{snippetexample}{elementary-matrix-five-by-five-example}{A row operation matrix}
    In \(\matrices_{5\times5}(\realnumbers)\),
    \[
        E_{45}(-7)
        =
        \begin{pmatrix}
            1 & 0 & 0 & 0 & 0\\
            0 & 1 & 0 & 0 & 0\\
            0 & 0 & 1 & 0 & 0\\
            0 & 0 & 0 & 1 & -7\\
            0 & 0 & 0 & 0 & 1
        \end{pmatrix}.
    \]
    Left multiplication by \(E_{45}(-7)\) performs the row operation
    \[
        R_4\leftarrow R_4-7R_5.
    \]
\end{snippetexample}

\begin{snippetproposition}{elementary-matrix-triangularity}{Triangular elementary matrices}
    Let \(\mathbb{F}\) be a \field, let \(n\in\naturalnumbers\), let
    \(1\leq i,j\leq n\) with \(i\neq j\), and let
    \(\lambda\in\mathbb{F}\). If \(i>j\), then \(E_{ij}(\lambda)\) is lower
    triangular. If \(i<j\), then \(E_{ij}(\lambda)\) is upper triangular.
\end{snippetproposition}

\begin{snippetproof}{elementary-matrix-triangularity-proof}{elementary-matrix-triangularity}{Triangular elementary matrices}
    The \matrix \(E_{ij}(\lambda)\) has the same entries as \(I_n\) except
    possibly in position \((i,j)\). If \(i>j\), then this exceptional entry
    lies below the main diagonal, so every entry above the main diagonal is
    zero and the \matrix is lower triangular. If \(i<j\), then the exceptional
    entry lies above the main diagonal, so every entry below the main diagonal
    is zero and the \matrix is upper triangular.
\end{snippetproof}

\section{Gaussian elimination in matrix form}

\begin{snippet}{first-elimination-step-matrix-form}
    Let \(\mathbb{F}\) be a \field, and let
    \[
        A=
        \begin{pmatrix}
            a_{11} & a_{12} & \cdots & a_{1n}\\
            a_{21} & a_{22} & \cdots & a_{2n}\\
            \vdots & \vdots & \ddots & \vdots\\
            a_{n1} & a_{n2} & \cdots & a_{nn}
        \end{pmatrix}
        \in\matrices_n(\mathbb{F}).
    \]
    If \(a_{11}\neq0\), the first pivot is \(a_{11}\). Eliminating the
    entries below it is obtained by the operations
    \[
        R_i\leftarrow R_i-\frac{a_{i1}}{a_{11}}R_1,
        \qquad
        i=2,\ldots,n.
    \]
    In \matrix form this first elimination step is
    \[
        E_{n1}\left(-\frac{a_{n1}}{a_{11}}\right)
        \cdots
        E_{31}\left(-\frac{a_{31}}{a_{11}}\right)
        E_{21}\left(-\frac{a_{21}}{a_{11}}\right)A,
    \]
    and it has the form
    \[
        \begin{pmatrix}
            a_{11} & a_{12} & \cdots & a_{1n}\\
            0 & a'_{22} & \cdots & a'_{2n}\\
            \vdots & \vdots & \ddots & \vdots\\
            0 & a'_{n2} & \cdots & a'_{nn}
        \end{pmatrix}.
    \]
\end{snippet}

\begin{snippetproposition}{triangular-matrix-product-proposition}{Product of triangular matrices}
    Let \(\mathbb{F}\) be a \field. The product of two upper triangular
    \matrix[matrices] in \(\matrices_n(\mathbb{F})\) is upper triangular. The product
    of two lower triangular \matrix[matrices] in \(\matrices_n(\mathbb{F})\) is lower
    triangular.
\end{snippetproposition}

\begin{snippetproof}{triangular-matrix-product-proposition-proof}{triangular-matrix-product-proposition}{Product of triangular matrices}
    Let \(A=(a_{ij})\) and \(B=(b_{ij})\) be lower triangular \matrix[matrices] in
    \(\matrices_n(\mathbb{F})\). We show that \(AB\) is lower triangular.
    If \(i<j\), then
    \[
        (AB)_{ij}=\sum_{k=1}^n a_{ik}b_{kj}.
    \]
    For each index \(k\), either \(k>i\), in which case \(a_{ik}=0\), or
    \(k\leq i<j\), in which case \(b_{kj}=0\). Hence every term in the sum
    is zero, so \((AB)_{ij}=0\). Thus \(AB\) is lower triangular. The proof
    for upper triangular \matrix[matrices] is analogous.
\end{snippetproof}

\begin{snippet}{gaussian-elimination-elementary-matrix-product}
    Let \(\mathbb{F}\) be a \field and let \(A\in\matrices_n(\mathbb{F})\).
    Gaussian elimination is a finite sequence of left multiplications by
    elementary \matrix[matrices]. If no row swaps are required and all pivots are
    nonzero, then
    \[
        \widetilde{E}_{n-1}\cdots\widetilde{E}_1A=U,
    \]
    where \(U\) is upper triangular and the \matrix[matrices]
    \(\widetilde{E}_k\) are lower triangular elementary \matrix[matrices]. If row
    swaps occur, elementary permutation \matrix[matrices] also appear:
    \[
        E_rP_r\cdots E_1P_1A=U.
    \]
\end{snippet}

\section{Invertibility}

\includesnpt{linearalgebra-matrix-inverse-definition}
\includesnpt{linearalgebra-invertible-matrix-definition}
\includesnpt{linearalgebra-uniqueness-of-matrix-inverse}
\includesnpt{linearalgebra-uniqueness-of-matrix-inverse-proof}

\begin{snippetproposition}{elementary-matrix-inverse}{Inverse of an elementary matrix}
    Let \(\mathbb{F}\) be a \field, let \(n\in\naturalnumbers\), let
    \(1\leq i,j\leq n\) with \(i\neq j\), and let
    \(\lambda\in\mathbb{F}\). The elementary \matrix \(E_{ij}(\lambda)\) is
    invertible and
    \[
        E_{ij}(\lambda)^{-1}=E_{ij}(-\lambda).
    \]
\end{snippetproposition}

\begin{snippetproof}{elementary-matrix-inverse-proof}{elementary-matrix-inverse}{Inverse of an elementary matrix}
    Left multiplication by \(E_{ij}(\lambda)\) performs the row operation
    \[
        R_i\leftarrow R_i+\lambda R_j.
    \]
    The inverse row operation is
    \[
        R_i\leftarrow R_i-\lambda R_j,
    \]
    which is left multiplication by \(E_{ij}(-\lambda)\). Hence
    \[
        E_{ij}(-\lambda)E_{ij}(\lambda)=I_n
        \qquad\text{and}\qquad
        E_{ij}(\lambda)E_{ij}(-\lambda)=I_n.
    \]
    Therefore \(E_{ij}(\lambda)^{-1}=E_{ij}(-\lambda)\).
\end{snippetproof}

\includesnpt{linearalgebra-product-rule-of-inverse-matrices}
\includesnpt{linearalgebra-product-rule-of-inverse-matrices-proof}
\includesnpt{linearalgebra-involution-rule-for-matrix-inverse}
\includesnpt{linearalgebra-involution-rule-for-matrix-inverse-proof}

\begin{snippetexample}{sum-of-invertible-matrices-warning}{A sum of invertible matrices need not be invertible}
    Let \(n\in\naturalnumbers\). The \matrix[matrices] \(I_n\) and \(-I_n\) are
    invertible, but
    \[
        I_n+(-I_n)=0,
    \]
    and the zero \matrix is not invertible.
\end{snippetexample}

\section{LU factorization}

\begin{snippettheorem}{lu-decomposition-theorem}{LU decomposition}
    Let \(\mathbb{F}\) be a \field and let
    \(A\in\matrices_n(\mathbb{F})\). Suppose that Gaussian elimination on
    \(A\) requires no row swaps and that all pivots are nonzero. Then
    \[
        A=LU,
    \]
    where \(L\in\matrices_n(\mathbb{F})\) is lower triangular with all
    diagonal entries equal to \(1\), and \(U\in\matrices_n(\mathbb{F})\) is
    upper triangular with the pivots of the elimination on its diagonal.
\end{snippettheorem}

\begin{snippetproof}{lu-decomposition-theorem-proof}{lu-decomposition-theorem}{LU decomposition}
    Since no row swaps are required, Gaussian elimination is obtained by left
    multiplying \(A\) by lower triangular elementary \matrix[matrices]
    \[
        \widetilde{E}_1,\ldots,\widetilde{E}_{n-1}
    \]
    with diagonal entries equal to \(1\). Thus
    \[
        \widetilde{E}_{n-1}\cdots\widetilde{E}_1A=U,
    \]
    where \(U\) is upper triangular and has the pivots on its diagonal. Each
    elementary \matrix is invertible, so
    \[
        A=(\widetilde{E}_{n-1}\cdots\widetilde{E}_1)^{-1}U.
    \]
    Put
    \[
        L=(\widetilde{E}_{n-1}\cdots\widetilde{E}_1)^{-1}.
    \]
    The inverse of each \(\widetilde{E}_k\) is again lower triangular with
    diagonal entries equal to \(1\), and products of lower triangular
    \matrix[matrices] are lower triangular. Hence \(L\) is lower triangular with all
    diagonal entries equal to \(1\), and \(A=LU\).
\end{snippetproof}

\begin{snippetproposition}{lu-decomposition-solving-systems}{Solving a system with LU decomposition}
    Let \(\mathbb{F}\) be a \field, let \(A\in\matrices_n(\mathbb{F})\), and
    let \(b\in\mathbb{F}^n\). If \(A=LU\), where \(L\) is lower triangular
    and \(U\) is upper triangular, then solving
    \[
        Ax=b
    \]
    is equivalent to solving the two triangular systems
    \[
        Ly=b,
        \qquad
        Ux=y.
    \]
\end{snippetproposition}

\begin{snippetproof}{lu-decomposition-solving-systems-proof}{lu-decomposition-solving-systems}{Solving a system with LU decomposition}
    Since \(A=LU\), the equation \(Ax=b\) is
    \[
        LUx=b.
    \]
    Put \(y=Ux\). Then \(LUx=b\) becomes \(Ly=b\). Conversely, if
    \(Ly=b\) and \(Ux=y\), then \(LUx=Ly=b\), so \(x\) solves \(Ax=b\).
\end{snippetproof}

\begin{snippetexample}{lu-factorization-example}{LU factorization}
    Consider
    \[
        A=
        \begin{pmatrix}
            1 & 4 & -2\\
            1 & -1 & 1\\
            4 & 0 & 2
        \end{pmatrix},
        \qquad
        b=
        \begin{pmatrix}
            -1\\
            -2\\
            -7
        \end{pmatrix}.
    \]
    The row operations
    \[
        R_2\leftarrow R_2-R_1,
        \qquad
        R_3\leftarrow R_3-4R_1
    \]
    give
    \[
        E_{31}(-4)E_{21}(-1)A
        =
        \begin{pmatrix}
            1 & 4 & -2\\
            0 & -5 & 3\\
            0 & -16 & 10
        \end{pmatrix}.
    \]
    Then
    \[
        R_3\leftarrow R_3-\frac{16}{5}R_2
    \]
    gives
    \[
        U=
        E_{32}\left(-\frac{16}{5}\right)E_{31}(-4)E_{21}(-1)A
        =
        \begin{pmatrix}
            1 & 4 & -2\\
            0 & -5 & 3\\
            0 & 0 & \frac{2}{5}
        \end{pmatrix}.
    \]
    Therefore
    \[
        L=
        \begin{pmatrix}
            1 & 0 & 0\\
            1 & 1 & 0\\
            4 & \frac{16}{5} & 1
        \end{pmatrix},
        \qquad
        A=LU.
    \]
    Solving \(Ly=b\) gives
    \[
        y=
        \begin{pmatrix}
            -1\\
            -1\\
            \frac{1}{5}
        \end{pmatrix}.
    \]
    Solving \(Ux=y\) gives
    \[
        x=
        \begin{pmatrix}
            -2\\
            \frac{1}{2}\\
            \frac{1}{2}
        \end{pmatrix}.
    \]
\end{snippetexample}

\section{Permutation matrices}

\begin{snippetdefinition}{elementary-permutation-matrix-definition}{Elementary permutation matrix}
    Let \(\mathbb{F}\) be a \field, let \(n\in\naturalnumbers\), and let
    \(1\leq i,j\leq n\) with \(i\neq j\). The \emph{elementary permutation
    \matrix} \(P_{ij}\in\matrices_n(\mathbb{F})\) is the \matrix obtained from
    \(I_n\) by swapping the \(i\)-th row with the \(j\)-th row.
\end{snippetdefinition}

\begin{snippetexample}{elementary-permutation-matrix-example}{An elementary permutation matrix}
    In \(\matrices_{4\times4}(\mathbb{F})\), the elementary permutation
    \matrix that swaps the second and fourth rows is
    \[
        P_{24}
        =
        \begin{pmatrix}
            1 & 0 & 0 & 0\\
            0 & 0 & 0 & 1\\
            0 & 0 & 1 & 0\\
            0 & 1 & 0 & 0
        \end{pmatrix}.
    \]
\end{snippetexample}

\begin{snippetproposition}{elementary-permutation-matrix-row-swap}{Effect of an elementary permutation matrix}
    Let \(\mathbb{F}\) be a \field, let \(A\in\matrices_{n\times m}(\mathbb{F})\),
    and let \(1\leq i,j\leq n\) with \(i\neq j\). Left multiplication by
    \(P_{ij}\) swaps the \(i\)-th and \(j\)-th rows of \(A\).
\end{snippetproposition}

\begin{snippetproof}{elementary-permutation-matrix-row-swap-proof}{elementary-permutation-matrix-row-swap}{Effect of an elementary permutation matrix}
    The identity \matrix satisfies \(I_nA=A\). The \matrix \(P_{ij}\) is
    obtained from \(I_n\) by swapping rows \(i\) and \(j\). Therefore
    \(P_{ij}A\) is obtained from \(I_nA=A\) by the same row swap.
\end{snippetproof}

\begin{snippetproposition}{elementary-permutation-matrix-inverse}{Inverse of an elementary permutation matrix}
    Let \(\mathbb{F}\) be a \field, let \(n\in\naturalnumbers\), and let
    \(1\leq i,j\leq n\) with \(i\neq j\). The elementary permutation \matrix
    \(P_{ij}\in\matrices_n(\mathbb{F})\) is invertible and
    \[
        P_{ij}^{-1}=P_{ij}.
    \]
\end{snippetproposition}

\begin{snippetproof}{elementary-permutation-matrix-inverse-proof}{elementary-permutation-matrix-inverse}{Inverse of an elementary permutation matrix}
    Swapping two rows twice returns every row to its original position. Hence
    left multiplication by \(P_{ij}P_{ij}\) has the same effect as left
    multiplication by \(I_n\). Therefore
    \[
        P_{ij}P_{ij}=I_n,
    \]
    and \(P_{ij}^{-1}=P_{ij}\).
\end{snippetproof}

\section{Gauss-Jordan method}

\begin{snippetproposition}{gauss-jordan-inverse-method}{Computing inverses with Gauss-Jordan}
    Let \(\mathbb{F}\) be a \field and let
    \(A\in\matrices_n(\mathbb{F})\). If a finite sequence of elementary row
    operations transforms \(A\) into \(I_n\), then applying the same
    operations to the augmented \matrix
    \[
        [A\mid I_n]
    \]
    gives
    \[
        [I_n\mid A^{-1}].
    \]
\end{snippetproposition}

\begin{snippetproof}{gauss-jordan-inverse-method-proof}{gauss-jordan-inverse-method}{Computing inverses with Gauss-Jordan}
    Let \(M\) be the product of the elementary \matrix[matrices] corresponding to the
    row operations. The hypothesis says that
    \[
        MA=I_n.
    \]
    Since elementary \matrix[matrices] are invertible, \(M\) is invertible. The
    equality \(MA=I_n\) implies that \(M=A^{-1}\). Applying the same row
    operations to \([A\mid I_n]\) left multiplies both blocks by \(M\), giving
    \[
        [MA\mid MI_n]=[I_n\mid M]=[I_n\mid A^{-1}].
    \]
\end{snippetproof}

\begin{snippetexample}{gauss-jordan-inverse-example}{Computing an inverse with Gauss-Jordan}
    Let
    \[
        A=
        \begin{pmatrix}
            1 & 4 & -2\\
            1 & -1 & 1\\
            4 & 0 & 2
        \end{pmatrix}.
    \]
    Starting from
    \[
        \left[
            \begin{array}{ccc|ccc}
                1 & 4 & -2 & 1 & 0 & 0\\
                1 & -1 & 1 & 0 & 1 & 0\\
                4 & 0 & 2 & 0 & 0 & 1
            \end{array}
        \right],
    \]
    perform
    \[
        R_2\leftarrow R_2-R_1,
        \qquad
        R_3\leftarrow R_3-4R_1,
        \qquad
        R_3\leftarrow R_3-\frac{16}{5}R_2.
    \]
    This gives
    \[
        \left[
            \begin{array}{ccc|ccc}
                1 & 4 & -2 & 1 & 0 & 0\\
                0 & -5 & 3 & -1 & 1 & 0\\
                0 & 0 & \frac{2}{5} & -\frac{4}{5} & -\frac{16}{5} & 1
            \end{array}
        \right].
    \]
    Normalize the second and third pivots:
    \[
        R_2\leftarrow -\frac{1}{5}R_2,
        \qquad
        R_3\leftarrow \frac{5}{2}R_3.
    \]
    Then eliminate above the pivots by
    \[
        R_2\leftarrow R_2+\frac{3}{5}R_3,
        \qquad
        R_1\leftarrow R_1+2R_3,
        \qquad
        R_1\leftarrow R_1-4R_2.
    \]
    The result is
    \[
        \left[
            \begin{array}{ccc|ccc}
                1 & 0 & 0 & 1 & 4 & -1\\
                0 & 1 & 0 & -1 & -5 & \frac{3}{2}\\
                0 & 0 & 1 & -2 & -8 & \frac{5}{2}
            \end{array}
        \right].
    \]
    Hence
    \[
        A^{-1}
        =
        \begin{pmatrix}
            1 & 4 & -1\\
            -1 & -5 & \frac{3}{2}\\
            -2 & -8 & \frac{5}{2}
        \end{pmatrix}.
    \]
\end{snippetexample}

\section{Nonsingularity}

\begin{snippetdefinition}{nonsingular-matrix-definition}{Nonsingular matrix}
    Let \(\mathbb{F}\) be a \field and let \(A\in\matrices_n(\mathbb{F})\).
    The \matrix \(A\) is \emph{nonsingular} if, for every
    \(b\in\mathbb{F}^n\), the linear system
    \[
        Ax=b
    \]
    has a unique solution.
\end{snippetdefinition}

\begin{snippettheorem}{nonsingular-matrix-invertible-theorem}{Nonsingularity and invertibility}
    Let \(\mathbb{F}\) be a \field and let \(A\in\matrices_n(\mathbb{F})\).
    Then \(A\) is nonsingular \ifandonlyif \(A\) is invertible.
\end{snippettheorem}

\begin{snippetproof}{nonsingular-matrix-invertible-theorem-proof}{nonsingular-matrix-invertible-theorem}{Nonsingularity and invertibility}
    Suppose first that \(A\) is invertible. For every \(b\in\mathbb{F}^n\),
    the vector
    \[
        x=A^{-1}b
    \]
    satisfies \(Ax=b\). If \(y\) is another solution, then
    \[
        y=A^{-1}Ay=A^{-1}b=x,
    \]
    so the solution is unique. Hence \(A\) is nonsingular.

    Conversely, suppose that \(A\) is nonsingular. Then Gaussian elimination
    produces a pivot in every column: if a pivot column were missing, the
    homogeneous system \(Ax=0\) would have a free variable and hence more than
    one solution. Continuing with Gauss-Jordan row operations transforms
    \(A\) into \(I_n\). By the Gauss-Jordan inverse method, these operations
    construct \(A^{-1}\). Hence \(A\) is invertible.
\end{snippetproof}

\end{document}
